\section*{第 \ref{ch:g8:sexual-reproduction} 章——有性生殖}

\begin{solution}{exo:g8:sexual-reproduction:1}
\omterm{def:g8:sexual-reproduction:gametes}{配子}是生殖\omterm{def:g5:first-look-at-cells:cell}{细胞}，有\omterm{def:g4:animal-reproduction:sexual}{雄性}和\omterm{def:g4:animal-reproduction:sexual}{雌性}之分；
\omterm{prop:g8:sexual-reproduction:core}{受精卵}是二者结合形成的单个新\omterm{def:g5:first-look-at-cells:cell}{细胞}。\omterm{def:g1:animals-around-us:animal}{动物}：
\omterm{def:g4:animal-reproduction:sexual}{精子}和\omterm{def:g4:animal-reproduction:sexual}{卵细胞}。被子\omterm{def:g1:plants-around-us:plant}{植物}：\omterm{def:g3:flowers-fruits-seeds:parts}{花粉}
粒中的\omterm{def:g4:animal-reproduction:sexual}{雄性}\omterm{def:g5:first-look-at-cells:cell}{细胞}与\omterm{prop:g5:flower-to-fruit:roles}{胚珠}中的\omterm{def:g4:animal-reproduction:sexual}{卵细胞}。
\end{solution}

\begin{solution}{exo:g8:sexual-reproduction:2}
两个亲本产生\omterm{def:g4:animal-reproduction:sexual}{雄性}\omterm{def:g8:sexual-reproduction:gametes}{配子}和\omterm{def:g4:animal-reproduction:sexual}{雌性}\omterm{def:g8:sexual-reproduction:gametes}{配子}；\omterm{def:g5:flower-to-fruit:steps}{受精}使二者各一个
结合成\omterm{prop:g8:sexual-reproduction:core}{受精卵}；\omterm{prop:g8:sexual-reproduction:core}{受精卵}不断分裂，
发育成一个新个体，它带有来自
双亲的遗传物质。
\end{solution}

\begin{solution}{exo:g8:sexual-reproduction:3}
\omterm{def:g8:sexual-reproduction:gametes}{配子}相遇的场所；发育中幼体的栖身之所；
数量与照料的策略。这些\omterm{def:g6:exploring-our-environment:conditions}{条件}
都围绕同一核心。
\end{solution}

\begin{solution}{exo:g8:sexual-reproduction:4}
\omterm{def:g3:flowers-fruits-seeds:parts}{雄蕊}：雄\omterm{def:g8:sexual-reproduction:gametes}{配子}的“工厂”。\omterm{def:g3:flowers-fruits-seeds:parts}{雌蕊}（及其\omterm{prop:g5:flower-to-fruit:roles}{胚珠}）：
雌\omterm{def:g8:sexual-reproduction:gametes}{配子}所在之处。\omterm{def:g3:flowers-fruits-seeds:parts}{花粉}管：体内\omterm{def:g5:flower-to-fruit:steps}{受精}的通道。
\omterm{def:g2:from-seed-to-plant:seed}{种子}：内含\omterm{def:g5:human-reproduction:pregnancy}{胚胎}及其营养储备的结构，是\omterm{def:g1:plants-around-us:plant}{植物}的
\omterm{def:g4:animal-reproduction:ovivivi}{卵生}“包裹”。
\end{solution}

\begin{solution}{exo:g8:sexual-reproduction:5}
单亲产生完全相同的后代——这种一致性是重组无法做到的；
双亲每次都产生新的组合——这种多样性是复制无法
提供的。
\end{solution}

\begin{solution}{exo:g8:sexual-reproduction:6}
每个同胞个体都由一对不同的\omterm{def:g8:sexual-reproduction:gametes}{配子}发育而来，
而且没有两个\omterm{def:g8:sexual-reproduction:gametes}{配子}携带的遗传信息完全相同，
每次\omterm{def:g5:flower-to-fruit:steps}{受精}都会把双亲的贡献重新组合，
产生新的组合。
\end{solution}

\begin{solution}{exo:g8:sexual-reproduction:7}
复制（嫁接、\omterm{met:g4:plants-without-seeds:cutting}{扦插}）保留了已被证明的优良性状，
适用于当前品种恰好符合需要的情形；其代价是缺陷也被共同继承。
重组（\omterm{def:g2:from-seed-to-plant:seed}{种子}繁殖）能探索新的组合，
适用于需要改良或保险的情形；
其代价是大多数组合并不出色。
\end{solution}

\begin{solution}{exo:g8:sexual-reproduction:8}
\omterm{def:g8:sexual-reproduction:gametes}{配子}：\omterm{def:g4:animal-reproduction:sexual}{雌性}产的卵和\omterm{def:g4:animal-reproduction:sexual}{雄性}产的\omterm{def:g4:animal-reproduction:sexual}{精子}，都排放到池塘中。
相遇：水中体外\omterm{def:g5:flower-to-fruit:steps}{受精}，\omterm{def:g4:animal-reproduction:sexual}{雄性}的鸣叫召集整个池塘。
栖身：卵包裹在胶质中，个体小且有营养储备。
数量与照料：数以千计，且无人照料。
季节：春季——五月的合唱。
\end{solution}

\begin{solution}{exo:g8:sexual-reproduction:9}
“相遇”一项不成立——根本没有\omterm{def:g5:flower-to-fruit:steps}{受精}——
而栖身、数量和季节三项都正常：夏季大量产出
带冠毛的\omterm{def:g2:from-seed-to-plant:seed}{种子}。这一例子说明，档案只是一份核对清单，
而不是既定假设：应依据证据来填写，
“\omterm{def:g2:from-seed-to-plant:seed}{种子}”不一定意味着重组。
（蒲公英以\omterm{def:g2:from-seed-to-plant:seed}{种子}进行无性繁殖，借助传播服务
进行复制式的“赌博”——两种情况兼而有之。）
\end{solution}

\begin{solution}{exo:g8:sexual-reproduction:10}
两百年的\omterm{met:g4:plants-without-seeds:cutting}{扦插}使葡萄园成为一个共同的遗传体——
优良性状被保留，缺陷也被共同继承。
喜爱这一遗传体的害虫发现每株葡萄对它都同样可口：
没有多样的个体，没有后备，复制这一“赌博”
在最不利的时刻被兑现。历史上的解决办法，
正如本题所示，是另一种“赌博”：
与抗性砧木杂交——即重组。
\end{solution}

\begin{solution}{exo:g8:sexual-reproduction:11}
从核心出发：单个\omterm{def:g8:sexual-reproduction:gametes}{配子}会死亡——只有\omterm{def:g5:flower-to-fruit:steps}{受精}，
即两个\omterm{def:g8:sexual-reproduction:gametes}{配子}结合，才能形成能够存活并分裂的\omterm{prop:g8:sexual-reproduction:core}{受精卵}。
从重组出发：每个\omterm{def:g8:sexual-reproduction:gametes}{配子}只携带一方亲本的一半组合，
这种组合直到结合时才会存在，
由\omterm{def:g5:flower-to-fruit:steps}{受精}形成。
\end{solution}

\begin{solution}{exo:g8:sexual-reproduction:12}
\omterm{ex:g4:plants-without-seeds:runner}{匍匐茎}是近距离、可靠的服务：在旁边优良的土地上密集地长出
完全相同的植株——这是当年已被证明的优良植株，
被保存下来。\omterm{def:g2:from-seed-to-plant:seed}{种子}则是远距离、探索性的服务：
重组形成的新组合被广泛散布，
为\omterm{def:g5:biodiversity:species}{物种}在别处变化的\omterm{def:g6:exploring-our-environment:conditions}{条件}下积累后备。
两者兼备，
就同时兼顾了当下与未来。
\end{solution}

\begin{solution}{pb:g8:sexual-reproduction:1}
\textbf{1.} 毛刷充当\omterm{def:g5:flower-to-fruit:steps}{传粉}者：将雄\omterm{def:g8:sexual-reproduction:gametes}{配子}从 A 植株
带到 B 植株的\omterm{def:g3:flowers-fruits-seeds:parts}{雌蕊}上——这是相遇的一步。套袋可防止蜜蜂
进入，使外来\omterm{def:g3:flowers-fruits-seeds:parts}{花粉}不会带来第三方亲本：这是对杂交的控制——
这是对杂交的控制，体现了公平实验的精神，
就在\omterm{def:g1:plants-around-us:parts}{花朵}上。

\textbf{2.} \omterm{def:g2:from-seed-to-plant:seed}{种子}——每粒\omterm{def:g2:from-seed-to-plant:seed}{种子}都是\omterm{def:g5:human-reproduction:pregnancy}{胚胎}，由 B 植株的一个\omterm{def:g4:animal-reproduction:sexual}{卵细
胞}与 A 植株的一个\omterm{def:g3:flowers-fruits-seeds:parts}{花粉}\omterm{def:g5:first-look-at-cells:cell}{细胞}\omterm{def:g5:flower-to-fruit:steps}{受精}而成：
结合了 A 与 B 两方的遗传物质。

\textbf{3.} 因为每粒\omterm{def:g2:from-seed-to-plant:seed}{种子}都是重组中的一次不同“组合”：
双方的\omterm{def:g8:sexual-reproduction:gametes}{配子}各不相同，因此会产生三百种不同的组合，
在脆度、抗性等方面各不相同，
其他方面也是如此。

\textbf{4.} 果园押下的两注：
（\cref{ex:g8:sexual-reproduction:bets}）：多数\omterm{prop:g2:from-seed-to-plant:cycle}{幼苗}平平无奇，
少数更差，极少数更好——
这正是洗牌已知的分散。

\textbf{5.} 这一组合又被洗牌打散：胜出者的
\omterm{def:g2:from-seed-to-plant:seed}{种子}与\omterm{def:g3:flowers-fruits-seeds:parts}{花粉}亲本的遗传物质组成\emph{新}的混合，
脆甜且抗性强的组合便分散为片段。
播种胜出者，反而会使它失去优势。

\textbf{6.} 嫁接——无性繁殖：亲本只有一个，
子代完全相同。它保证了珍贵的组合，
一棵接一棵地复制。

\textbf{7.} 复制式赌注的共同弱点风险：
全园只有同一种体质——葡萄园的故事
迟早会重演。

\textbf{8.} 多样性储备：其他品种间种，
并从\omterm{def:g2:from-seed-to-plant:seed}{种子}保存育种株系——
一旦出现新的病害或气候要求，洗牌便有了可用的材料。
这是\omterm{def:g5:biodiversity:species}{物种}层面的储备，
在果园尺度上有意保留。

\textbf{9.} 第一阶段（杂交与三百株\omterm{prop:g2:from-seed-to-plant:cycle}{幼苗}）：
洗牌——它\emph{发现}了这一组合。
第二阶段（嫁接胜出者）：复制——
它\emph{保留}了这一组合。先探索，
后保存。

\textbf{10.} 差异之处：相遇（体外，在水中，无需人工授粉——育种者只需配对亲本），
产卵（卵产于巢中，而非\omterm{def:g2:from-seed-to-plant:seed}{种子}），照料（雄鱼扇动卵块），
季节。不变之处是：产生\omterm{def:g8:sexual-reproduction:gametes}{配子}、
\omterm{def:g5:flower-to-fruit:steps}{受精}使两者结合为合子，合子发育为
经过重组的个体——
这一核心不能改变，
否则便不是\omterm{def:g4:animal-reproduction:sexual}{有性生殖}。

\textbf{11.} 洗牌不断发牌——每一代都发出无人设计的组合；
育种者的挑选，
只是有意识地执行保留原则：
自然\omterm{def:g6:exploring-our-environment:conditions}{条件}变化时挑选\omterm{def:g4:adapted-to-their-world:adaptation}{适应}者，
育种者则挑选脆甜且抗性强的。
没有发出的品种，便无可挑选。

\textbf{12.} \omterm{def:g4:animal-reproduction:sexual}{有性生殖}是为了\emph{产生}新的组合；
无性繁殖则是为了
\emph{原样地增殖}它。
\end{solution}
